% =============================================================================
% Project : THE_PHILOSOPHY_OF_VIABILITY
% Book    : Book I — Ontology of Viability
% Book ID : VS.B01
% Part    : Part I — Identity and Manifestation (VS.B01.P01)
% Chapter : Chapter 08 — Comparative Philosophical Examination
% Chapter ID : VS.B01.P01.C08
%
% File    : 4_ASSETS/diagrams/VS.B01.P01.C08.S05.SS07.P005.DG01.potentiality_actuality_and_form.tex
%
% Section / Chapter / Asset:
%   Diagram DG01 — Potentiality, Actuality, and Form
%
% Structural Identifier : VS.B01.P01.C08.S05.SS07.P005.DG01
% Asset Type : DG
%
% Purpose :
%   To display the explanatory relation among matter, potentiality, form, and
%   actuality without representing them as a universal linear hierarchy of
%   being.
%
% Representational limit :
%   The arrangement is explanatory rather than a universal sequence or
%   hierarchy of all beings.
%
% Dependencies :
%   TikZ; libraries used: arrows.meta, positioning. All are loaded by the
%   project preamble; no per-file package or library loading.
%
% =============================================================================

\begin{figure}[htbp]
    \centering
    \hyphenpenalty=10000 \exhyphenpenalty=10000\relax  % no hyphenation inside diagram labels

    \begin{tikzpicture}[
        font=\small,
        % --- base box ---
        box/.style={
            draw, rounded corners=3pt, align=center,
            text width=2.9cm, minimum height=1.5cm,
            inner sep=5pt, line width=0.7pt
        },
        % --- node roles ---
        concept/.style={box, fill=blue!8, draw=blue!55!black},
        glossbox/.style={box, dashed, fill=gray!10, draw=black!60},
        % --- arrows and labels ---
        directed/.style={-{Stealth[length=2.5mm]}, line width=0.9pt},
        vslink/.style={draw=black!55, dashed, line width=0.7pt},
        elabel/.style={font=\scriptsize\itshape, text=black!75, align=center}
    ]

    % ---------------------------------------------------------------
    % Explanatory arrangement: matter - potentiality - form - actuality
    % ---------------------------------------------------------------
    \node[concept] (matter) {
        \textbf{Matter}\\[2pt] \textit{hylē}
    };
    \node[concept, right=2.6cm of matter] (potentiality) {
        \textbf{Potentiality}\\[2pt] \textit{dynamis}
    };
    \node[concept, below=2.0cm of potentiality] (form) {
        \textbf{Form}\\[2pt] \textit{eidos / morphē}
    };
    \node[concept, right=2.6cm of form] (actuality) {
        \textbf{Actuality}\\[2pt] \textit{energeia / entelecheia}
    };

    % ---------------------------------------------------------------
    % Glosses (dashed, undirected)
    % ---------------------------------------------------------------
    \node[glossbox, below=1.2cm of matter] (matternote) {
        Underlying subject and capacity for determination
    };
    \node[glossbox, below=1.2cm of actuality] (actualitynote) {
        Realization and determinate constitution
    };
    \draw[vslink] (matter)    -- (matternote);
    \draw[vslink] (actuality) -- (actualitynote);

    % ---------------------------------------------------------------
    % Relations
    % ---------------------------------------------------------------
    \draw[directed] (matter) -- node[elabel, above=2pt, text width=2.4cm]
        {considered as capable} (potentiality);
    \draw[directed] (potentiality) -- node[elabel, right=4pt, text width=2.6cm]
        {relation to determination} (form);
    \draw[directed] (form) -- node[elabel, above=2pt, text width=2.4cm]
        {actualizing role} (actuality);

    \end{tikzpicture}

    \caption{A schematic presentation of the relation between matter,
    potentiality, form, and actuality in Aristotle's hylomorphic
    analysis. The arrangement is explanatory rather than a
    universal sequence or hierarchy of all beings.}
    \label{fig:aristotle-potentiality-actuality-form}
\end{figure}

% =============================================================================
% END OF DIAGRAM DG01
% =============================================================================
